\itemhead{4.2 思考 · 未编号题组（教材p.106）}{XJ-B1-C04-TH-P106-Y0720}
\label{XJ-B1-C04-TH-P106-Y0720}
{\footnotesize 湘教版教材p.106 / PDF第113页 / 4.2 思考 / 未编号题组}
\field{原题}{\subq{1}将一张纸对折一次得到2层，对折两次得到4层。对折 $x$ 次后，纸的层数 $y$ 与 $x$ 之间有什么关系？
\subq{2}《庄子》中说：“一尺之棰，日取其半，万世不竭。”若每天截取木棰的一半，剩余木棰的长度 $y$（尺）与截取的天数 $x$ 之间有什么关系？
\subq{3}上述两个关系有什么共同特征？}
\field{参考答案}{（1）$y=2^x$，$x\in\mathbb N$；（2）$y=(1/2)^x$，$x\in\mathbb N$，表示每天截取后剩余长度；（3）均为等时间间隔按固定正倍率变化的指数模型，离散情境中自变量取非负整数。}
{\small \field{关联知识}{E03 指数增长、衰减与实际模型（教材帮p.131、132、154 / PDF140、141、163）}
\field{方法与易错点}{确认初始量、时间单位、倍率与时刻起点，再代入或逐次比较阈值。 易错：把年折旧10\%写成乘0.1；“超过”用成“达到”；时间起点错一位；把衰减至接近0当实际为0。}
\field{答案依据}{答案册未提供本栏目的独立对应答案；依据原题、教材概念补写参考回应。开放性总结不设唯一数值答案。}}
